% Cote 119: offprints, reports, courses, CV and letters (typed and printed part): archive cote 119, site page 76 (PDF page 77).
% offprint. Draft compared: none.
% AI-transcribed by Claude Opus 5.5 (claude-opus-5-5, adaptive thinking, effort high), Message Batches API, batch r119_full (msgbatch_01DwZNm5Y7YsKBWyAN4sFs8w), 2026-11-13; not yet checked by eye.
\scanpage{119/76}
\mspage{( 4 )}

On déduit de 2.3 et 2.4 que si $N$ est une sous-catégorie tannakienne de $M_k$
possédant un $\otimes$-générateur, il existe des théories de la cohomologie définies
sur $N$ et à valeurs dans une extension finie de $\mathbf{Q}$.

3.2. Soit $k$ un corps parfait de caractéristique $p \neq 0$, W l'anneau des
vecteurs de Witt sur $k$, K son corps des fractions; l'automorphisme de
Frobenius de $k$ induit un automorphisme $\sigma$ du corps K. On appelle F-iso-%
cristal sur $k$ un espace vectoriel M de rang fini sur K, munie d'un automor-%
phisme $\sigma$-linéaire $\mathrm{F}_{\mathrm{M}} : \mathrm{M} \to \mathrm{M}$, i. e., vérifiant $\mathrm{F}_{\mathrm{M}}(\lambda x) = \sigma(\lambda)\, \mathrm{F}_{\mathrm{M}}(x)$ si
$\lambda \in \mathrm{K}$, $x \in \mathrm{M}$. On obtient ainsi une catégorie F Criso$(k)$, qui est muni d'une
loi $\otimes$ ACU évidente, et pour laquelle c'est une catégorie tannakienne sur
le corps $\mathbf{Q}_p$ des nombres $p$-adiques.

La catégorie tannakienne F Criso$(k)$ dépend fonctoriellement du corps
parfait $k$. Si $k$ est algébriquement clos, on prouve que la gerbe tannakienne
sur $\mathbf{Q}_p$ qui la définit est liée par un $\mathbf{Q}_p$-groupe pro-diagonalisable, ayant
$\mathbf{Q}$ comme groupe de caractères.

Les F-isocristaux apparaissent dans la théorie de Dieudonné-Grothendieck
des groupes de Barsotti-Tate, et en cohomologie cristalline [(\textsuperscript{1}), (\textsuperscript{5})]. On
espère que la cohomologie cristalline des variétés lisses et projectives sur $k$
définit un morphisme de catégories tannakiennes
\[ \mathrm{M}_k \to \mathrm{F\,Cris}\,(k). \]

(*) Séance du 25 janvier 1971.
(\textsuperscript{1}) P. Berthelot, Comptes rendus, 269, série A, 1969, p. 297.
(\textsuperscript{2}) M. Demazure, Motifs des variétés algébriques, Séminaire Bourbaki, 365
(novembre 1969).
(\textsuperscript{3}) P. Gabriel, Bull. Soc. Math. Fr., 90, 1962, p. 323-448.
(\textsuperscript{4}) J. Giraud, Cohomologie non-abélienne, Notes miméographiées, Columbia University.
(\textsuperscript{5}) A. Grothendieck, TDTE III, Séminaire Bourbaki, 195 (février 1969).
(\textsuperscript{6}) A. Grothendieck, Notes by J. Coates and O. Jussila : Crystals and the De Rham
cohomology of schemes; dix exposés sur la cohomologie des schémas, North-Holland, 1968.
(\textsuperscript{7}) A. Grothendieck, Standard Conjectures on Algebraic Cycles; Proceedings of the
Bombay Colloquium on Algebraic Geometry, 1968, p. 193-199.
(\textsuperscript{8}) S. Kleiman, Algebraic cycles and the Weil conjectures; dix exposés sur la cohomologie
des schémas, North-Holland, 1968.
(\textsuperscript{9}) N. Saavedra, Comptes rendus, 272, série A, 1971, p. 258.

(Institut
des Hautes Études scientifiques,
35, route de Chartres,
91-Bures-sur-Yvette, Essonne.)

183281. --- Imp. Gauthier-Villars. --- 55, Quai des Grands-Augustins, Paris (6\textsuperscript{e}).
Imprimé en France.

\mspage{\hand{76}}
